\section{Compact transverse-traceless propagation and the combined bound}
\label{sec:ttcombined}

\subsection{General compact quadrupole radiation}

Let $Q_{ij}$ be an arbitrary possibly complex symmetric trace-free quadrupole amplitude at angular frequency $\omega$. For propagation direction $\hat{\bm n}$ define
\begin{equation}
P_{ij}=\delta_{ij}-n_i n_j
\end{equation}
and the transverse-traceless projector
\begin{equation}
\Lambda_{ij,kl}(\hat{\bm n})
=\frac12\left(
P_{ik}P_{jl}+P_{il}P_{jk}-P_{ij}P_{kl}
\right).
\label{eq:TTprojector}
\end{equation}
For unit-normalized polarization tensors, completeness gives
\begin{equation}
\sum_\lambda
|Q_{ij}\epsilon_{ij}^{(\lambda)}|^2
=Q_{ij}^*\Lambda_{ij,kl}Q_{kl}.
\end{equation}
In the quantum normalization the angular one-graviton rate is
\begin{equation}
\frac{\dd\kappa_g}{\dd\Omega}
=
\frac{G\omega^5}{4\pi\hbar c^5}
Q^*:\Lambda(\hat{\bm n}):Q,
\label{eq:differentialrate}
\end{equation}
but the normalized angular dependence below is identical for a classical monochromatic TT field. Rotational invariance and the fact that $\Lambda$ has rank two on the five-dimensional STF tensor space imply
\begin{equation}
\int\dd\Omega\,\Lambda(\hat{\bm n})
=\frac{8\pi}{5}I_{\rm STF}.
\label{eq:LambdaIntegral}
\end{equation}
Thus
\begin{equation}
\kappa_g
=
\frac{2G\omega^5}{5\hbar c^5}Q^*:Q,
\end{equation}
consistent with Eq.~\eqref{eq:gravitonrate}.

The corresponding normalized directivity is
\begin{equation}
D_Q(\hat{\bm n})
=\frac52\frac{Q^*:\Lambda(\hat{\bm n}):Q}{Q^*:Q}.
\end{equation}
Because $\Lambda$ is an orthogonal projector,
\begin{equation}
\boxed{D_Q(\hat{\bm n})\le\frac52.}
\label{eq:Dmax}
\end{equation}
Equality occurs when $Q$ lies entirely in the two-dimensional TT tensor subspace associated with the line of sight. For real STF tensors this directivity functional is algebraically identical to the compact gravitational-antenna result of Hirakawa, Narihara, and Fujimoto \cite{HirakawaNariharaFujimoto1976}; neither the functional nor the value $5/2$ is claimed as new.

\subsection{Normalized TT propagation singular value}

A convenient tensor-normalized angular radiation mode is
\begin{equation}
 u_{Q,\lambda}(\hat{\bm n})
=
\sqrt{\frac{5}{8\pi}}
\frac{Q:\epsilon^{(\lambda)}(\hat{\bm n})}{\sqrt{Q^*:Q}},
\label{eq:normalizedmode}
\end{equation}
which obeys
\begin{equation}
\sum_\lambda\int\dd\Omega\,
|u_{Q,\lambda}|^2=1.
\end{equation}
The same normalized mode may be read as a one-graviton angular wavefunction or as an energy-normalized classical TT radiation pattern. For normalized source and receiver quadrupoles separated by $R\hat{\bm R}$, fixed-frequency translation gives
\begin{equation}
S_{BA}(z)
=
\sum_\lambda\int\dd\Omega\,
 u_{Q_B,\lambda}^*(\hat{\bm n})
 u_{Q_A,\lambda}(\hat{\bm n})
\ee^{iz\hat{\bm n}\cdot\hat{\bm R}},
\qquad z=kR.
\label{eq:translatedoverlap}
\end{equation}
In the wave zone, the outgoing stationary point at $+\hat{\bm R}$ gives
\begin{equation}
\boxed{
t_{BA}^{\rm TT}(z)
=
-\frac{5i}{4z}\ee^{iz}
\frac{Q_B^*:\Lambda(\hat{\bm R}):Q_A}
{\sqrt{Q_A^*:Q_A}\sqrt{Q_B^*:Q_B}}
+O(z^{-2}).}
\label{eq:TTamplitude}
\end{equation}
Cauchy--Schwarz and $\|\Lambda\|_{\op}=1$ therefore imply
\begin{equation}
\boxed{
\|P_g(\omega)\|_{\op}^2
\le
\frac{25}{16[k(\omega)R]^2}
\left[1+O((kR)^{-1})\right].}
\label{eq:TTpropbound}
\end{equation}
Matched source and receiver tensors in the same line-of-sight TT subspace saturate the leading projector inequality. For the aligned plus tensor $Q\propto\mathrm{diag}(1,-1,0)$, the complete normalized angular integral reduces to
\begin{equation}
S(z)=\frac5{32}
\int_{-1}^{1}\dd x\,
(1+6x^2+x^4)\ee^{izx},
\end{equation}
whose outgoing endpoint contribution is $-(5i/4z)\ee^{iz}+O(z^{-2})$. This independently fixes the coefficient in Eq.~\eqref{eq:TTamplitude}.

The leading coefficient also has a purely classical reciprocal-antenna interpretation after the endpoint coupling magnitudes have been normalized out. Two matched compact quadrupoles have maximal directivities $D_A=D_B=5/2$, and the far-field normalized transfer factor is
\begin{equation}
D_AD_B\left(\frac{\lambda}{4\pi R}\right)^2
=\frac{25}{16(kR)^2},
\label{eq:FriisTT}
\end{equation}
using $\lambda=2\pi/k$. Equation~\eqref{eq:TTamplitude} is the direct TT-field normalization of the same leading propagation geometry; quantum mechanics is not required for the coefficient.

\subsection{Passive compact-quadrupole throughput theorem}

We now combine the passive endpoint cut set, the classical cumulative quadrupole oscillator-strength bound, and normalized TT propagation. For a narrow band centered at $\omega$, assume
\begin{equation}
\frac{\omega L_A}{c}\ll1,
\qquad
\frac{\omega L_B}{c}\ll1,
\qquad
kR=\frac{\omega R}{c}\gg1.
\label{eq:regime}
\end{equation}
Then Eqs.~\eqref{eq:materialcutset} and \eqref{eq:TTpropbound} give the main result.

\begin{theorem}[Passive compact-quadrupole gravitational throughput]
For direct weak one-way propagation between compact passive nonrelativistic linear endpoints in the regime \eqref{eq:regime},
\begin{equation}
\boxed{
\Gamma_{\rm coh}
\lesssim
\frac{25}{16(kR)^2}
\frac{4G\omega^4}{3c^5}
\min(I_A,I_B).}
\label{eq:combined1}
\end{equation}
Equivalently,
\begin{equation}
\boxed{
\Gamma_{\rm coh}
\lesssim
\frac{25G\omega^2}{12c^3R^2}
\min(I_A,I_B).}
\label{eq:combined2}
\end{equation}
\end{theorem}

The headline inequality is a classical passive-wave bound within the stated linear-harmonic class: the H2 cut set, cumulative modal resource ceiling, and leading propagation bound all admit classical derivations. Quantization reproduces the same endpoint normalization and supplies the information-theoretic corollaries discussed later. The inequality is leading-order in the compact-source and wave-zone expansions and is an upper bound rather than a claim that all constituent inequalities can be saturated simultaneously.

Its main structural feature is what no longer appears: endpoint quality factor, branching fraction, number of passive internal resonances, and the basis used to describe coherent internal mode mixing. Those quantities determine how closely a particular device approaches the ceiling, but they cannot increase the ceiling inside the stated class.

For a finite band $\mathcal B=[\omega_-,\omega_+]$ lying entirely in the compact wave-zone regime, a conservative broad-band version keeps the two band edges distinct:
\begin{equation}
\Gamma_{\rm coh}(\mathcal B)
\lesssim
\frac{25G}{12c^3R^2}
\frac{\omega_+^4}{\omega_-^2}
\min(I_A,I_B).
\label{eq:broadband}
\end{equation}
This form is intentionally loose; a sharper broadband response theory should retain the full frequency dependence inside the integral rather than combining separate suprema.
